\chapter{हमारी ज़िम्मेदारी}\label{ch:g9:our-responsibility}

इस किताब के नौ साल यहीं ख़त्म होते हैं, और उसके सारे धागे एक सीधे तथ्य पर
आ मिलते हैं: जो \omterm{def:g5:biodiversity:species}{प्रजाति} यह पन्ना पढ़ सकती है वही अकेली \omterm{def:g5:biodiversity:species}{प्रजाति} है जो उस
मशीनरी को समझती है जिसका वह हिस्सा है --- और समझ, जैसा
\cref{ch:g8:contraception} ने पहली बार कहा था, चुनाव लाती है। यह आख़िरी
अध्याय उन्हीं ज़िम्मेदारियों को समेटता है जो यह समझ पैदा करती है: अपनी
सेहत की ओर, दूसरों की सेहत की ओर, और सजीवों के उस जाल की ओर, जो हम सबको
ढो रहा है।

\section{अपनी ही मशीनरी की ज़िम्मेदारी}

\begin{proposition}[सेहत, यानी जानकारी वाली अपनी सँभाल]\label{prop:g9:our-responsibility:self}
इस किताब के सेहत वाले धागे, बँधे हुए:
\begin{itemize}
  \item रोज़ का इक़रारनामा
        (\cref{prop:g5:healthy-choices:contract}) --- खाना, हरकत,
        \omterm{prop:g1:healthy-habits:needs}{नींद}, सफ़ाई --- अंगों का वह रखरखाव है जिसे अब तुम तंत्र-दर-तंत्र
        सही ठहरा सकते हो
        (\cref{prop:g7:heart-and-circulation:keep},
        \cref{met:g7:lungs-and-gas-exchange:keep},
        \cref{met:g8:protecting-the-nervous-system:manual});
  \item उन पदार्थों की कार्य-प्रणालियाँ मालूम हैं
        (\cref{prop:g8:protecting-the-nervous-system:substances}), और
        उस जाल के बेमेल की क़ीमत भी लगी हुई है
        (\cref{rem:g5:healthy-choices:never});
  \item जनन वाले चुनावों का अपना नक़्शा और अपने पते हैं
        (\cref{ch:g8:contraception});
  \item और शरीर का इतिहास कुछ हद तक बाँटा हुआ है
        (\cref{ch:g9:heredity-and-traits}): किसी परिवार के रुझान जानना
        --- जल्दी सफ़ेद होते बाल, परत जमते \omterm{def:g4:heart-and-blood:heart}{दिल} --- कोई नियति नहीं, एक
        रखरखाव की समय-सारणी है
        (\cref{met:g9:unique-individuals:claims}: सीमाएँ, जो जी जाती
        हैं)।
\end{itemize}
\end{proposition}
\admitted

\begin{example}[जाँच, जीव विज्ञान की तरह पढ़ी हुई]\label{ex:g9:our-responsibility:checkup}
कोई आम जाँच एक मुलाक़ात में सिमटी यही किताब है: \omterm{prop:g4:heart-and-blood:pulse}{नाड़ी} और दबाव (सड़क के
किनारे पढ़ा गया वह जाल,
\cref{ex:g7:heart-and-circulation:measures}), प्रतिवर्त वाला हथौड़ा
(\cref{pb:g8:nervous-communication:1}), \omterm{prop:g9:immune-defenses:vaccination}{टीकाकरण} की बही
(\cref{pb:g9:immune-defenses:1}), परिवार का इतिहास (यानी बाँटे हुए
रुझान), और वही पक्की सलाह --- यानी फिर वह इक़रारनामा। दवा बड़ी हद तक उन्हीं
कार्य-प्रणालियों का व्यावहारिक रखरखाव है, जिनकी पुस्तिकाएँ अब तुम्हारे
पास हैं।
\end{example}

\section{दूसरों की ओर ज़िम्मेदारी}

\begin{proposition}[सेहत एक साझी मशीनरी है]\label{prop:g9:our-responsibility:others}
तीन कार्य-प्रणालियाँ सेहत को एक साझा खाता बना देती हैं:
\begin{itemize}
  \item \emph{संचरण}
        (\cref{def:g9:microbes-and-infection:stages}): तुम्हारी
        \omterm{def:g1:healthy-habits:hygiene}{स्वच्छता}, तुम्हारी खाँसी का ढका जाना और तुम्हारा बीमारी वाले दिन
        घर रहना, ये सब दूसरों की हिफ़ाज़त हैं;
  \item \emph{साझी \omterm{def:g9:immune-defenses:memory}{प्रतिरक्षा}}
        (\cref{prop:g9:immune-defenses:vaccination}): तुम्हारा टीका
        तुम्हारे ज़रिए उनकी रखवाली करता है जिन्हें सिखाया ही नहीं जा
        सकता;
  \item \emph{दवा की साझी पूँजी}: चलते हुए बने रहे प्रतिजैविक
        (\cref{exo:g9:how-species-change:10}) और दिया गया \omterm{prop:g4:journey-of-food:absorb}{रक्त}
        (\cref{ex:g9:unique-individuals:medicine} --- वह बही दानियों पर
        ही चलती है) ऐसी साझी चीज़ें हैं जिन्हें अलग-अलग लोगों के चुनाव
        ख़र्च भी करते हैं और भरते भी।
\end{itemize}
\end{proposition}
\admitted

\begin{example}[एक किशोर का साझा खाता]\label{ex:g9:our-responsibility:account}
उस हफ़्ते की आम प्रविष्टियाँ: बुख़ार में घर रहा (संचरण कटा); टीके ताज़ा
थे (वह साझापन थमा रहा); वसंत का प्रतिजैविक कोर्स पूरा किया (साझी पूँजी
बिना ख़र्च की); और अठारह पर पहला रक्तदान तय (वह बही भरी हुई)। इनमें कुछ
भी वीरता नहीं --- यानी \cref{ex:g5:healthy-choices:day} वाला सबक़, पूरे
क़स्बे के पैमाने पर: जन-स्वास्थ्य आम दिन ही हैं, सँभाले हुए।
\end{example}

\section{सजीव दुनिया की ओर ज़िम्मेदारी}

\begin{proposition}[रखवाली, अपने सारे कारणों के साथ दोबारा कही हुई]\label{prop:g9:our-responsibility:stewardship}
प्रकृति वाले धागे, अपनी-अपनी कार्य-प्रणालियों के साथ बँधे हुए:
\begin{itemize}
  \item हर \omterm{def:g3:food-chains:chain}{आहार शृंखला} हरे से शुरू होती है
        (\cref{prop:g6:plants-are-producers:firstlink}), और हर थाली
        किसी परितंत्र की पैदावार है
        (\cref{prop:g6:producing-our-food:costs});
  \item जाल झटके विविधता से सोखते हैं
        (\cref{prop:g5:biodiversity:web}), और विविधता ही वह बचत है
        जिस पर चयन चलता है
        (\cref{prop:g9:unique-individuals:reserve},
        \cref{exo:g9:how-species-change:12}) --- यानी \omterm{def:g5:biodiversity:biodiversity}{जैव विविधता} कोई
        नज़ारा नहीं, चलती हुई पूँजी है;
  \item विलुप्ति हमेशा के लिए दरवाज़ा बंद कर देती है
        (\cref{prop:g5:biodiversity:loss}) --- और वह इकलौता
        जान-बूझकर किया गया अपवाद, चेचक, साबित करता है कि यह शब्द कितना
        अंतिम है (\cref{exo:g9:immune-defenses:12});
  \item टिकाऊ इस्तेमाल के नियम मालूम हैं
        (\cref{prop:g5:people-and-nature:lasting}), उनके उल्लंघन
        पहचाने जा सकते हैं
        (\cref{ex:g7:oxygen-in-the-water:waste}), और उनकी मरम्मतें कर
        के दिखाई जा चुकी हैं
        (\cref{ex:g5:people-and-nature:river},
        \cref{ex:g5:biodiversity:storks})।
\end{itemize}
रखवाली वही रखरखाव वाला मिज़ाज है जो सेहत का है, बस उस बड़े शरीर पर लागू
किया हुआ जिसके भीतर हम रहते हैं।
\end{proposition}
\admitted

\begin{example}[किसी नागरिक का जीव विज्ञान किसलिए है]\label{ex:g9:our-responsibility:citizen}
बड़ों की ज़िंदगी के वोट, ख़रीदारियाँ और आदतें बार-बार ऐसे सवालों के भेस
में आती हैं जिनके लिए यह किताब तैयार कर चुकी है: क्या यह मछली-उद्योग
भरपाई से तेज़ ले रहा है? क्या यह नीति परागकर्ताओं का \omterm{def:g2:where-animals-live:habitat}{आवास} बचाती है या उस
पर सड़क बना देती है? क्या यह चमत्कारी परहेज़, इलाज या ``फ़लाँ का \omterm{def:g9:genes-and-dna:gene}{जीन}''
वाला दावा तुक में बैठता है
(\cref{met:g9:unique-individuals:claims},
\cref{met:g8:contraception:evaluate})? क्या इस नाँद में यह प्रतिजैविक
होना चाहिए? जो नागरिक ये जाँचें चला सकता है, नौ साल के अध्याय उसी के लिए
थे।
\end{example}

\begin{method}[स्नातक का तरीक़ा]\label{met:g9:our-responsibility:method}
पहली किताब छोड़ते हुए, तथ्यों से ज़्यादा यह तरीक़ा साथ ले जाओ:
\begin{enumerate}
  \item नतीजा निकालने से पहले देखो
        (\cref{met:g1:animals-around-us:observe} ने यह शुरू किया था;
        और उसके बाद की हर गिनती तथा हर नोटबुक ने);
  \item निष्पक्ष परखो --- एक फ़र्क़, और एक \omterm{met:g4:what-plants-need:fairtest}{नियंत्रण}
        (\cref{met:g4:what-plants-need:fairtest}, जो उसके बाद हर
        कक्षा में इस्तेमाल हुआ);
  \item कार्य-प्रणाली का नाम लो --- जब तक किसी दावे की कड़ी का नाम न
        बताया जाए तब तक उसकी जाँच हुई ही नहीं
        (\cref{met:g8:contraception:evaluate});
  \item आज़ाद गवाहों को तौलो
        (\cref{met:g9:evidence-of-evolution:weigh});
  \item और \omterm{def:g6:exploring-our-environment:components}{सजीव} मशीनरी का लिहाज़ करो, अपने शरीर में भी और उसके बाहर भी:
        वह पुरानी है, बारीक है, हदों के भीतर मरम्मत लायक़ है --- और
        सँभालना तुम्हें ही है।
\end{enumerate}
\end{method}

\begin{remark}[आगे क्या आता है]\label{rem:g9:our-responsibility:next}
हाई स्कूल वाला खंड इन्हीं धागों को अगले आवर्धन पर उठाता है: \omterm{def:g5:first-look-at-cells:cell}{कोशिका} की
आणविक मशीनरी, अपने पूरे गणित के साथ आनुवंशिकी, \omterm{def:g9:immune-defenses:memory}{प्रतिरक्षा} तथा \omterm{def:g5:brain-in-command:system}{तंत्रिका
तंत्र} संजीदगी से, पारिस्थितिकी एक संख्यात्मक विज्ञान की तरह, और \omterm{prop:g9:how-species-change:selection}{जैव विकास}
उस ढाँचे के रूप में जो वह है। सवाल वही रहेंगे, बस और पास से पूछे हुए।
तरीक़ा साथ लाना।
\end{remark}

\section{अभ्यास}

\begin{exercise}[$\star$]\label{exo:g9:our-responsibility:1}
\cref{prop:g9:our-responsibility:self} में अपनी ज़िम्मेदारी के चारों
धागे गिनाओ।
\end{exercise}

\begin{exercise}[$\star$]\label{exo:g9:our-responsibility:2}
सेहत को साझा बनाने वाली तीनों कार्य-प्रणालियों के नाम बताओ, और हर एक के
साथ उसका अध्याय भी।
\end{exercise}

\begin{exercise}[$\star$]\label{exo:g9:our-responsibility:3}
परिवार का जाना-पहचाना रुझान नियति के बजाय समय-सारणी क्यों है?
\end{exercise}

\begin{exercise}[$\star$]\label{exo:g9:our-responsibility:4}
जाँच के तीन औज़ारों को उनके अध्यायों की तरह पढ़ो।
\end{exercise}

\begin{exercise}[$\star$]\label{exo:g9:our-responsibility:5}
\omterm{def:g5:biodiversity:biodiversity}{जैव विविधता} नज़ारे के बजाय ``चलती हुई पूँजी'' क्यों है? दो
कार्य-प्रणालियाँ बताओ।
\end{exercise}

\begin{exercise}[$\star$]\label{exo:g9:our-responsibility:6}
स्नातक के तरीक़े के पाँचों क़दम याददाश्त से बताओ।
\end{exercise}

\begin{exercise}[$\star\star$]\label{exo:g9:our-responsibility:7}
बीमारी वाले दिन घर रहने को जन-स्वास्थ्य का काम साबित करो, और चरण का नाम
भी लो।
\end{exercise}

\begin{exercise}[$\star\star$]\label{exo:g9:our-responsibility:8}
तुम्हारा टीका किसकी ढाल है? जवाब उस साझेपन से दो, और दो टिके हुए लोगों
के नाम भी बताओ।
\end{exercise}

\begin{exercise}[$\star\star$]\label{exo:g9:our-responsibility:9}
\cref{ex:g9:our-responsibility:citizen} की तरह जाँच करो: ``हमारा नया
छिड़काव बाग़ के हर \omterm{prop:g3:sorting-living-things:groups}{कीट} को मार देता है --- यानी पूरी हिफ़ाज़त।''
\end{exercise}

\begin{exercise}[$\star\star$]\label{exo:g9:our-responsibility:10}
``प्रतिजैविक एक साझी पूँजी हैं'' को उस वार्ड के चयन से समझाओ --- और वे
दोनों आदतें भी बताओ जो उस पूँजी को भरी रखती हैं।
\end{exercise}

\begin{exercise}[$\star\star$]\label{exo:g9:our-responsibility:11}
किसी वायरल ``चमत्कारी इलाज'' वाली पोस्ट पर उस तरीक़े के पाँचों क़दम, एक
के बाद एक, लागू करो।
\end{exercise}

\begin{exercise}[$\star\star\star$]\label{exo:g9:our-responsibility:12}
इस किताब का समापन अनुच्छेद ख़ुद लिखो: \omterm{def:g5:first-look-at-cells:cell}{कोशिका}, शरीर, जाल और वंश-वृक्ष के
लिए एक-एक वाक्य --- और इन चारों में पढ़ने वाले की जगह भी।
\end{exercise}

\section{समस्या: स्नातक वाली जाँच-पड़ताल}

\begin{problem}[{सप्ताहांत समस्या --- किसी क़स्बे का एक साल, पूरी किताब
से जाँचा हुआ}]\label{pb:g9:our-responsibility:1}
तुम्हारे क़स्बे का (कल्पित) साल, सुर्ख़ियों में: सर्दी में फ़्लू की एक
लहर; वसंत में चिनार के \omterm{def:g1:plants-around-us:tree}{पेड़ों} पर छिड़काव को लेकर बहस; गर्मी में पार्क के
तालाब में \omterm{prop:g3:sorting-living-things:groups}{मछलियों} की मौत; पतझड़ में नई दलदल-शरणस्थली के लिए अभियान; और
साल भर, दवाख़ाने के \omterm{prop:g9:immune-defenses:vaccination}{टीकाकरण} तथा दान वाले अभियान।

\medskip\noindent\textbf{भाग I --- सेहत वाली फ़ाइल।}
\begin{enumerate}
  \item फ़्लू की लहर: बताओ कि क़स्बे की सलाह ने किन संचरण-रास्तों पर
        निशाना साधा, और दवाख़ाने ने सबसे पहले देखभाल-गृहों के
        कर्मचारियों को टीका क्यों लगाया।
  \item वे अभियान: बताओ कि \omterm{prop:g4:journey-of-food:absorb}{रक्त} की बही और \omterm{prop:g9:immune-defenses:vaccination}{टीकाकरण} की बही, हर एक क्या
        भरती है, और उन पर हक़ किसका है।
  \item एक स्तंभकार ``हाथ धोने पर हो-हल्ले'' का मज़ाक़ उड़ाता है।
        \cref{ex:g9:microbes-and-infection:history} वाले वार्ड से एक
        अनुच्छेद में जवाब दो।
  \item दवाख़ाना बताता है कि अधूरे छोड़े गए प्रतिजैविक कोर्स लौट रहे
        हैं। कौन-सी साझी पूँजी ख़र्च हो रही है, और किस कार्य-प्रणाली
        से?
\end{enumerate}

\medskip\noindent\textbf{भाग II --- प्रकृति वाली फ़ाइल।}
\begin{enumerate}[resume]
  \item छिड़काव वाली बहस: ``हर \omterm{prop:g3:sorting-living-things:groups}{कीट} मार दो'' पर बाग़ वाली जाँच चलाओ
        (\cref{exo:g9:our-responsibility:9}) --- और कौन-कौन मरता है,
        और कौन-सी मुफ़्त सेवाएँ रुक जाती हैं?
  \item \omterm{prop:g3:sorting-living-things:groups}{मछलियों} की मौत: एक गरम, ठहरा हुआ हफ़्ता और मटर जैसा हरा
        तालाब। \cref{ch:g7:oxygen-in-the-water} से रोग पहचानो, और
        नुस्ख़ा लिखो।
  \item दलदल वाला अभियान: उस शरणस्थली का बचाव तीन कार्य-प्रणालियों से
        करो --- \omterm{def:g2:where-animals-live:habitat}{आवास}, जाल, और विविधता की बचत --- हर एक पर एक वाक्य।
  \item एक ज़मींदार पूछते हैं, ``इस दलदल ने आख़िर क़स्बे के लिए किया
        क्या है?'' जवाब मुफ़्त काम की उस बही से दो
        (\cref{exo:g5:biodiversity:11}, और
        \cref{ex:g5:people-and-nature:river} वाला नदी का बाढ़-भंडार)।
\end{enumerate}

\medskip\noindent\textbf{भाग III --- स्नातक का समापन।}
\begin{enumerate}[resume]
  \item टिकाऊ इस्तेमाल के नियमों से उस साल का सबसे अच्छा और सबसे बुरा
        फ़ैसला चुनो, और हर फ़ैसले का औचित्य भी दो।
  \item दिखाओ कि \omterm{prop:g3:sorting-living-things:groups}{मछलियों} की मौत सुलझाने में उस तरीक़े के पाँचों क़दम
        कैसे काम आए --- पहले अवलोकन से लेकर मानी हुई कार्य-प्रणाली
        तक।
  \item क़स्बे के अख़बार का साल के आख़िर वाला संपादकीय चार वाक्यों में
        लिखो: सेहत साझी, प्रकृति बैंक में, दावे जाँचे हुए, और तरीक़ा
        सँभाला हुआ।
  \item पहली किताब के आख़िरी पन्ने की आख़िरी पंक्ति तुम्हें लिखनी है:
        एक वाक्य, उस पाठक के नाम, जिसने \cref{ch:g1:living-or-not}
        वाले घोंघे से शुरुआत की थी।
\end{enumerate}
\end{problem}
